\section{Dos conferencias, un método de investigación: primero mirar los datos, después estudiar el cambio}

Esta clase se compone de dos conferencias de estilo distinto y lógica complementaria. Jason Wei parte de los datos de entrenamiento y de las curvas de comportamiento para preguntarse por qué un next-token objective tan simple produce tantas capacidades; Hyung Won Chung, en cambio, parte de los cambios históricos del Transformer para preguntarse qué estructuras son atajos en las condiciones de cómputo actuales y cuáles se convertirán en cuellos de botella cuando la escala crezca. Al poner juntas las dos conferencias, no se obtienen dos grupos de opiniones dispersas, sino un método común: \textbf{primero entrar en contacto con evidencia real, después descomponer la observación en mecanismos verificables y, por último, estudiar las tendencias a lo largo del eje de la escala o del tiempo.}

\subsection{La pregunta de Jason: ¿por qué los modelos de lenguaje funcionan tan bien?}

Esta sección establece primero la postura de investigación de la primera conferencia. La portada no promete una teoría completa, sino que dice explícitamente que ofrecerá algunas intuitions. Aquí, una intuition no es una conjetura arbitraria, sino una hipótesis de trabajo condensada a partir de los datos, de las curvas de pérdida y de prompts concretos; debe poder orientar el siguiente experimento, y no solo explicar el éxito a posteriori.

\lecturefigure{jason-slide-001.jpg}{Tema de la conferencia de Jason Wei: Intuitions on Language Models}{Jason Wei official deck, p.~1}

\begin{importantbox}{Las cuatro preguntas de la primera conferencia}
\begin{enumerate}
  \item ¿Por qué la next-word prediction no es solo una tarea, sino algo parecido a una mezcla de una enorme cantidad de tareas?
  \item ¿Por qué aumentar la escala reduce de forma estable la loss global?
  \item Si la loss global es suave, ¿por qué una tarea individual presenta cambios abruptos?
  \item ¿Por qué ciertas tareas cuidadosamente seleccionadas muestran inverse o U-shaped scaling?
\end{enumerate}
\end{importantbox}

\lecturefigure{jason-slide-002.jpg}{La pregunta básica y el método de inspeccionar manualmente los datos de entrenamiento}{Jason Wei official deck, p.~2}

Al leer esta diapositiva, conviene mirar primero el método de la esquina inferior derecha, y no solo el título. ``Why do LLMs work?'' es demasiado amplia para experimentar con ella directamente; ``manually inspect data'', en cambio, convierte la pregunta en una observación ejecutable: tomar muestras aleatorias de texto, hacer la continuation a mano, marcar el conocimiento y las reglas que requiere un token y preguntarse después de qué contextos pudo haber aprendido el modelo esas capacidades. Esto vuelve a conectar la investigación abstracta sobre modelos con la distribución de entrenamiento.

\teachervoice{Entre 00:00:55 y 00:01:27, Jason llama a «mirar los datos a mano» una herramienta extremadamente útil. Esta advertencia se opone a mirar solo las puntuaciones agregadas de los benchmarks: si el investigador nunca ha leído las muestras de entrenamiento, es fácil que confunda una señal de supervisión evidente en los datos con una emergencia misteriosa.}

\lecturefigure{jason-slide-003.jpg}{Mirar los datos a mano es entrenar la propia «red neuronal biológica» del investigador}{Jason Wei official deck, p.~3}

El sentido didáctico de esta frase es que el investigador también necesita representation learning. Leer muestras una y otra vez, predecir etiquetas y revisar los errores con expertos del dominio forma en la mente representaciones del ruido, la ambigüedad, los atajos y la cola larga. Jason usa su experiencia con la clasificación de cáncer de pulmón para mostrar que aprender primero la tarea misma le dio la intuición que más tarde le permitió plantear preguntas de investigación; esto no exige que cada investigador se convierta en un experto completo del dominio, pero sí que entienda, al menos, qué significa un error.

\begin{lstlisting}[caption={Ciclo mínimo de investigación de la manual data inspection}]
for example in sample(training_corpus):
    guess = human_predict(example.context)
    compare(guess, example.next_token)
    label_required_signal(example)
    record_noise_shortcut_ambiguity(example)
summarize_recurring_patterns()
design_one_testable_experiment()
\end{lstlisting}

\subsection{La pregunta de Hyung Won: cuando no se alcanza a seguir los avances más recientes, estudiar el cambio mismo}

La subsección anterior construyó la intuición a partir de los datos; ahora se amplía la escala temporal. El título de Hyung Won subraya ``from the history of Transformer'', pero su propósito no es hacer una crónica, sino identificar en la historia las fuerzas impulsoras que se repiten. Al leer las siguientes diapositivas, ``history'' debe entenderse como un conjunto de experimentos naturales: cuando cambian el cómputo, los datos, los objetivos y la estructura, qué métodos se conservan, cuáles se descartan y por qué.

\lecturefigure{hyung-slide-001.jpg}{Tema de la conferencia de Hyung Won Chung: moldear el futuro de la IA a partir de la historia del Transformer}{Hyung Won Chung official deck, p.~1}

\lecturefigure{hyung-slide-002.jpg}{No perseguir solo los resultados más recientes: estudiar el cambio mismo}{Hyung Won Chung official deck, p.~2}

La segunda diapositiva redefine la ansiedad de investigación como un problema de método. Perseguir cada semana los modelos nuevos proporciona una gran cantidad de hechos, pero no necesariamente una dirección; estudiar el ``change itself'' exige, en cambio, alinear los sistemas antiguos con los nuevos y encontrar las variables comunes detrás de los cambios de desempeño. El objetivo no es predecir con precisión el nombre de algún producto, sino aumentar la probabilidad de acertar al juzgar la dirección de la investigación.

\teachervoice{Entre 00:32:29 y 00:33:22, Hyung Won enfatiza que, cuando la velocidad del cambio supera la capacidad humana de seguirlo, conviene más volver a mirar lo antiguo y trazar la ruta de «cómo llegamos hasta aquí». Lo importante de la clase no es memorizar la tabla de modelos de 2024, sino practicar la abstracción de una dirección a partir de la ruta.}

\lecturefigure{hyung-slide-003.jpg}{Método de tres pasos para estudiar el cambio: identificar, comprender y extrapolar la fuerza dominante}{Hyung Won Chung official deck, p.~3}

\begin{importantbox}{Definición de trabajo de dominant driving force}
La fuerza impulsora dominante no es la única fuerza del sistema, sino la que explica la mayor parte de la direccionalidad a la escala del problema actual. El investigador puede ignorar temporalmente los factores menores a cambio de poder analizar el sistema; pero debe declarar abiertamente que se trata de una aproximación y volver a examinar los factores ignorados cuando la predicción falle.
\end{importantbox}

\lecturefigure{hyung-slide-004.jpg}{Experimento de la pluma que cae: la gravedad para ilustrar cómo ignorar las fuerzas secundarias}{Hyung Won Chung official deck, p.~4}

Para una pluma que se suelta desde el reposo, si se toma la dirección hacia abajo como positiva, cerca de la superficie y despreciando la resistencia del aire, se tiene
\[
y(t)=y_0+v_0t+\frac{1}{2}gt^2.
\]
Aquí, $y(t)$ representa la posición en el instante $t$, $y_0$ es la posición inicial, $v_0$ es la velocidad inicial y $g$ es la aceleración de la gravedad. La fórmula es útil no porque la resistencia no exista, sino porque, con la precisión que requiere esta demostración, la gravedad domina lo suficiente y ya se cuenta con un modelo confiable de ella.

\teachervoice{Entre 00:35:04 y 00:36:21, la clase menciona explícitamente la resistencia del aire y las interacciones aerodinámicas complejas, pero decide ignorarlas. Lo que enseña este ejemplo no es que «la realidad tiene una sola variable», sino «primero definir la precisión requerida y después decidir qué variables pueden aproximarse».}

\lecturefigure{hyung-slide-006.jpg}{Cuantas más fuerzas dominantes y más complejas sus interacciones, más difícil es predecir la trayectoria}{Hyung Won Chung official deck, p.~6}

\begin{knowledgebox}{Lectura de la figura: no confundir el eje horizontal con el «número total de variables»}
El eje horizontal es el \textbf{número de fuerzas dominantes}, no el número de campos observados. Un sistema puede tener miles de características y, aun así, estar gobernado por unos pocos mecanismos; a la inversa, un sistema en apariencia simple también puede ser difícil de extrapolar si tiene varias retroalimentaciones fuertemente acopladas. Esta figura respalda una estrategia de modelado: primero buscar los factores con mayor direccionalidad y después decidir si hace falta incorporar sus interacciones.
\end{knowledgebox}

\subsection{El ciclo de investigación común}

Las dos subsecciones anteriores partieron, respectivamente, de un corte transversal de los datos y de un eje histórico longitudinal; ahora pueden combinarse en un ciclo. Primero se miran los datos reales para no alejarse de la distribución; después se trazan curvas o trayectorias históricas para identificar la dirección; luego se propone una explicación mínima del mecanismo y se buscan activamente contraejemplos que la refuten. La diferencia entre los dos conferencistas está solo en la escala de observación, no en la estructura epistemológica.

\begin{table}[H]
\centering
\begin{tabularx}{\textwidth}{L{0.19\textwidth} >{\raggedright\arraybackslash}X >{\raggedright\arraybackslash}X}
\toprule
Paso & Versión de Jason & Versión de Hyung Won \\
\midrule
Contacto con la evidencia & Leer a mano muestras de entrenamiento y prompts & Comparar arquitecturas antiguas con las actuales \\
Búsqueda de estructura & Descomponer latent tasks y loss mixture & Identificar la dominant driving force \\
Construcción de curvas & Trazar scaling curves de compute/data/performance & Extrapolar la fuerza impulsora a lo largo del tiempo y la escala \\
Revisión de límites & Examinar la métrica, la calidad de los datos y los puntos intermedios faltantes & Examinar las fuerzas ignoradas, los cuellos de botella físicos y el regime change \\
Paso a la acción & Diseñar el siguiente grupo de experimentos a escala & Agregar o retirar sesgos inductivos \\
\bottomrule
\end{tabularx}
\caption{El ciclo de investigación guiado por evidencia que comparten las dos conferencias.}
\end{table}

\subsection{Resumen de la sección}

Esta sección establece la forma de leer toda la clase. Jason pide al investigador que entrene su intuición sobre los datos; Hyung Won, que entrene su intuición histórica; ambos se oponen a saltar de un resultado aislado a conclusiones grandiosas. El contenido técnico que sigue repetirá una y otra vez la misma operación: descomponer el objetivo, entender las curvas, identificar los supuestos y explicar qué no puede concluirse a partir de la evidencia.
